\documentclass[letterpaper,journal]{IEEEtran}

\usepackage{amsmath,amssymb,amsfonts,bm}
\usepackage{array}
\usepackage{booktabs}
\usepackage{cite}
\usepackage{graphicx}
\usepackage{makecell}
\usepackage{multirow}
\usepackage{placeins}
\usepackage{url}
\usepackage{kotex}

\hyphenation{op-tical net-works semi-conduc-tor IEEE-Xplore}

\newcommand{\RMS}{\operatorname{RMS}}
\newcommand{\WRMS}{\operatorname{WRMS}}
\newcommand{\RWRMS}{\operatorname{RWRMS}}
\newcommand{\MTVV}{\operatorname{MTVV}}
\newcommand{\VDV}{\operatorname{VDV}}
\newcommand{\CF}{\operatorname{CF}}
\newcommand{\PtP}{\operatorname{P2P}}
\newcommand{\AbsPeak}{\operatorname{AbsPeak}}

\begin{document}

\title{Physics-Informed Candidate Features for Short-Window Ride Events}
\author{}
\maketitle

\section{Purpose and scope}

This document defines a \emph{candidate feature space}; it does not prescribe one universally
optimal KPI per sensor channel. The ride-analysis review in \cite{guastadisegni2024} explicitly
states that KPI suitability depends on the road profile, manoeuvre, measured quantity, and
subjective attribute. In particular, isolated obstacles are commonly described by peak and
peak-to-peak quantities, long road profiles by RMS-based quantities, and strongly transient
vibration by running-RMS or fourth-power quantities.

ISO~2631-1 and BS~6841 define ways to quantify human exposure to whole-body vibration
\cite{iso2631,bs6841}. They provide physically meaningful frequency weightings, but they do
not establish which feature best predicts the binary feedback in this data set. The ISO-based
quantities below are therefore used as physics-informed candidates. Their predictive relevance
must be determined using held-out-driver posterior-predictive evaluation.

\section{Measured and derived quantities}

Let $x(t)$ denote a signal over a window of duration $T$ and sampling rate $f_s$. The basic
descriptors used by the feature pipeline are
\begin{align*}
\RMS(x) &= \sqrt{\frac{1}{T}\int_0^T x^2(t)\,\mathrm{d}t},\\
\AbsPeak(x) &= \max_t |x(t)|,\\
\PtP(x) &= \max_t x(t)-\min_t x(t),\\
\CF_{\mathrm{raw}}(x) &= \frac{\AbsPeak(x)}{\RMS(x)}.
\end{align*}
The code names are \texttt{rms}, \texttt{abs\_peak}, \texttt{p2p}, and \texttt{crest},
respectively. These are generic signal descriptors. In particular,
$\CF_{\mathrm{raw}}$ is not the standards-based weighted crest factor defined below.

For a measured rate signal, differentiation produces the quantity used by the review's
acceleration-based KPIs:
\[
\dot{x}(t)=\frac{\mathrm{d}x(t)}{\mathrm{d}t}.
\]
Thus, differentiating pitch, roll, or yaw rate gives angular acceleration, while differentiating
heave velocity gives vertical acceleration. The corresponding code names are
\texttt{rms\_deriv}, \texttt{p2p\_deriv}, \texttt{abs\_peak\_deriv}, and
\texttt{vdv\_deriv}. Differentiation amplifies sensor noise, so smoothing, sampling rate, and
finite-difference accuracy must be treated as part of the feature definition.

\section{ISO frequency-weighted quantities}

For a translational or rotational acceleration $a(t)$, define
\[
a_w(t)=\mathcal{F}^{-1}\!\left\{W(f)\,\mathcal{F}\{a\}(f)\right\}.
\]
With $s=j2\pi f$ and $\omega_i=2\pi f_i$, the ISO weighting is constructed from
\begin{align*}
H_h(s) &= \frac{s^2}{s^2+\sqrt{2}\,\omega_1s+\omega_1^2},\\
H_\ell(s) &= \frac{\omega_2^2}{s^2+\sqrt{2}\,\omega_2s+\omega_2^2},\\
H_t(s) &= \frac{1+s/\omega_3}
{1+s/(Q_4\omega_4)+(s/\omega_4)^2},\\
H_s(s) &= \left(\frac{\omega_5}{\omega_6}\right)^2
\frac{1+s/(Q_5\omega_5)+(s/\omega_5)^2}
{1+s/(Q_6\omega_6)+(s/\omega_6)^2}.
\end{align*}
The scale-step normalization in $H_s$ is $({\omega_5}/{\omega_6})^2$. Reversing this
ratio preserves the frequency-response shape but multiplies $W_k$ by approximately four,
so it is not the ISO-normalized physical quantity.

The relevant parameter sets are summarized in Table~\ref{tab:weights}.
\begin{table}[t]
\centering
\caption{ISO~2631-1 weighting parameters used by the current channels.}
\label{tab:weights}
\resizebox{\columnwidth}{!}{%
\begin{tabular}{@{}lccccl@{}}
\toprule
Weight & $f_1$ & $f_2$ & $f_3=f_4$ & $Q_4$ & Scale step \\
\midrule
$W_k$ & 0.4 & 100 & 12.5 & 0.63 & $f_5=2.37$, $f_6=3.35$, $Q_5=Q_6=0.91$ \\
$W_d$ & 0.4 & 100 & 2.0  & 0.63 & none \\
$W_e$ & 0.4 & 100 & 1.0  & 0.63 & none \\
\bottomrule
\end{tabular}}
\end{table}
The resulting magnitudes are
\begin{align*}
W_k &= |H_hH_\ell H_tH_s|,\\
W_d &= |H_hH_\ell H_t|,\\
W_e &= |H_hH_\ell H_t|.
\end{align*}
where the transition parameters differ as shown in Table~\ref{tab:weights}. Their intended
directions are
\begin{itemize}
\item $W_k$: vertical translational acceleration ($z$ direction),
\item $W_d$: longitudinal and lateral translational acceleration ($x,y$ directions),
\item $W_e$: rotational acceleration about the roll, pitch, and yaw axes.
\end{itemize}

The standards-based weighted RMS and transient quantities are
\begin{align*}
\WRMS(a;W) &= \sqrt{\frac{1}{T}\int_0^T a_w^2(t)\,\mathrm{d}t},\\
\CF_w(a;W) &= \frac{\max_t|a_w(t)|}{\WRMS(a;W)},\\
\RWRMS(t_0) &= \sqrt{\frac{1}{\tau_0}
\int_{t_0-\tau_0}^{t_0} a_w^2(t)\,\mathrm{d}t},\\
\MTVV(a;W) &= \max_{t_0}\RWRMS(t_0),\\
\VDV(a;W) &= \left(\int_0^T |a_w(t)|^4\,\mathrm{d}t\right)^{1/4}.
\end{align*}
The review uses $\CF_w=6$ as a BS~6841-based stationary/transient guideline, whereas
ISO~2631-1 normally regards the basic weighted-RMS method as sufficient up to a crest
factor of 9. Neither threshold should be applied to the current unweighted
\texttt{crest} feature.

\subsection{Current implementation status}

\begin{table}[t]
\centering
\caption{Interpretation of existing feature-pipeline names.}
\label{tab:implementation}
\begin{tabular}{@{}ll@{}}
\toprule
Code name & Interpretation \\
\midrule
\texttt{wrms\_z}, \texttt{wrms\_xy} & frequency-weighted translational RMS \\
\texttt{wrms\_rot}, \texttt{wrms\_z\_d} & differentiate rate, then weighted RMS \\
\texttt{crest} & unweighted generic peak-to-RMS ratio \\
\texttt{vdv}, \texttt{vdv\_deriv} & unweighted fourth-power feature \\
\texttt{mtvv} & unweighted running-RMS maximum \\
\bottomrule
\end{tabular}
\end{table}
The last three rows are useful exploratory statistics, but they must not be reported as ISO
$\CF$, $\VDV$, or $\MTVV$ unless weighting is applied before their calculation.

\section{Signal-to-candidate mapping}

Table~\ref{tab:mapping} maps each available signal to physically meaningful candidates.
The ordering distinguishes event-oriented base candidates from standards-inspired or
exploratory additions; it is not a claim of predictive importance.

\begin{table*}[t]
\centering
\caption{Candidate metrics for the available channels. No row defines a universal primary KPI.}
\label{tab:mapping}
\resizebox{\textwidth}{!}{%
\begin{tabular}{@{}lllll@{}}
\toprule
Signal & Physical quantity & Event-oriented candidates & Additional candidates & Applicability / caveat \\
\midrule
\texttt{IMU\_LongAccelVal}
& longitudinal body acceleration $a_x$
& \texttt{abs\_peak}, \texttt{p2p}, \texttt{rms}
& \texttt{wrms\_xy}, \texttt{p95\_abs}, \texttt{vdv}, \texttt{crest}
& impact and traction/braking response; $W_d$ is directionally valid \\

\texttt{IMU\_VerAccelVal}
& vertical body acceleration $a_z$
& \texttt{abs\_peak}, \texttt{p2p}, \texttt{rms}
& \texttt{wrms\_z}, \texttt{p95\_abs}, \texttt{vdv}, \texttt{crest}, \texttt{mtvv}
& central ride signal; \texttt{z\_nod} is not valid on acceleration \\

\texttt{Pitch\_rate\_6D}
& pitch rate $\dot{\theta}$; derive $\ddot{\theta}$
& \texttt{abs\_peak\_deriv}, \texttt{p2p\_deriv}, \texttt{rms\_deriv}
& \texttt{wrms\_rot}, \texttt{vdv\_deriv}, raw-rate \texttt{rms}
& RMS/P2P angular acceleration is supported for pitch control and braking \\

\texttt{Roll\_rate\_6D}
& roll rate $\dot{\phi}$; derive $\ddot{\phi}$
& \texttt{abs\_peak\_deriv}, \texttt{p2p\_deriv}, \texttt{rms\_deriv}
& \texttt{wrms\_rot}, \texttt{vdv\_deriv}, raw-rate \texttt{rms}
& relevant to asymmetric roads, road copying, and cornering \\

\texttt{IMU\_YawRtVal}
& yaw rate $\dot{\psi}$; derive $\ddot{\psi}$
& \texttt{p2p\_deriv}, \texttt{rms\_deriv}
& \texttt{wrms\_rot}, raw-rate \texttt{rms}, \texttt{crest}
& mathematically valid rotation; mainly a handling/context candidate \\

\texttt{Bounce\_rate\_6D}
& heave velocity $\dot{z}$; derive $\ddot{z}$
& \texttt{abs\_peak\_deriv}, \texttt{p2p\_deriv}, \texttt{rms\_deriv}
& \texttt{wrms\_z\_d}, \texttt{vdv\_deriv}, raw-rate \texttt{rms}
& derived vertical acceleration is valid; displacement-decay KPIs are not \\

\texttt{IMU\_LatAccelVal}
& lateral body acceleration $a_y$
& \texttt{abs\_peak}, \texttt{p2p}, \texttt{rms}
& \texttt{wrms\_xy}, \texttt{p95\_abs}, \texttt{vdv}, \texttt{crest}
& most relevant to cornering or asymmetric events; $W_d$ is directionally valid \\

\texttt{IMU\_RollRtVal}
& IMU roll rate $\dot{\phi}$; derive $\ddot{\phi}$
& same as \texttt{Roll\_rate\_6D}
& use as an alternative sensor/source
& likely redundant with \texttt{Roll\_rate\_6D}; compare rather than assume both \\
\bottomrule
\end{tabular}}
\end{table*}

\section{Selection by event regime}

\begin{table*}[t]
\centering
\caption{The event regime determines which statistics are defensible candidates.}
\label{tab:regime}
\resizebox{\textwidth}{!}{%
\begin{tabular}{@{}lll@{}}
\toprule
Regime / target attribute & First candidates & Conditional additions \\
\midrule
Isolated bump, cleat, or short impact
& absolute peak, P2P, RMS of linear/angular acceleration
& weighted RMS; weighted RWRMS/MTVV/VDV for strong transients \\
Long, flat, rough, or approximately stationary road
& RMS and frequency-weighted RMS
& band-limited energy, vibration total value \\
Acceleration or braking pitch
& RMS/P2P pitch angular acceleration, longitudinal acceleration
& pitch settling/response time when a steady state is identifiable \\
Cornering or asymmetric road input
& RMS/P2P roll angular acceleration, lateral acceleration
& yaw and steering/context quantities \\
Oscillation decay after a step
& displacement settling time, overshoot, third displacement peak
& only when displacement and a defensible steady state are available \\
Road holding
& suspension deflection, suspension working space, dynamic wheel load
& unavailable from the current selected IMU channels alone \\
\bottomrule
\end{tabular}}
\end{table*}

For the current short windows around a reported event, the isolated-event row is the most
appropriate starting point. WRMS remains a useful perceptually weighted candidate, but it
should compete with peak, P2P, and RMS features rather than replace them. A four-second
window also provides limited low-frequency resolution; the resulting WRMS is better described
as an ISO-inspired relative feature than as a standards-compliant exposure measurement.

\FloatBarrier
\section{Evidence-based priorities for the current bump task}

Table~\ref{tab:bump-priority} translates the review's event- and attribute-specific evidence
into priorities for the available signals. ``First'' means that the review directly supports the
same physical response for an isolated impact or a closely related body-control attribute.
``Conditional'' means that the statistic is supported only for a particular event type, requires
an additional diagnostic, or is a physically reasonable extrapolation without a direct pairing in
Table~4 of \cite{guastadisegni2024}.

\begin{table*}[t]
\centering
\caption{Evidence-based candidate priority for the current short-window bump task.}
\label{tab:bump-priority}
\begingroup
\setlength{\tabcolsep}{3pt}
\small
\begin{tabular}{@{}
>{\raggedright\arraybackslash}p{0.18\textwidth}
>{\raggedright\arraybackslash}p{0.205\textwidth}
>{\raggedright\arraybackslash}p{0.245\textwidth}
>{\raggedright\arraybackslash}p{0.315\textwidth}@{}}
\toprule
Signal & First candidates & Conditional candidates & Evidence and rationale \\
\midrule
\texttt{IMU\_LongAccelVal}
& \texttt{p2p}
& \texttt{abs\_peak}, \texttt{rms}, \texttt{wrms\_xy}
& Figure~11 reports a direct relationship between longitudinal seat-rail P2P and bump impact harshness. The current CG-level IMU is a proxy for that measurement location. \\

\texttt{IMU\_VerAccelVal}
& \texttt{p2p}, \texttt{abs\_peak}, \texttt{rms}
& \texttt{wrms\_z}, \texttt{p95\_abs}, properly weighted \texttt{MTVV}/\texttt{VDV}
& Section~2.2 identifies vertical peak and P2P for isolated obstacles. Table~4 pairs bump impact with seat vertical P2P and pairs vertical/bounce control with vertical RMS acceleration. \\

\texttt{Pitch\_rate\_6D}
& \texttt{rms\_deriv}, \texttt{p2p\_deriv}
& \texttt{abs\_peak\_deriv}, \texttt{wrms\_rot}, settling/response time
& Table~4 pairs RMS pitch angular acceleration with pitch control and P2P/RMS pitch angular acceleration with pitch under acceleration or braking. Differentiation maps the measured rate to this quantity. \\

\texttt{Roll\_rate\_6D},
\texttt{IMU\_RollRtVal}
& \texttt{rms\_deriv}
& \texttt{p2p\_deriv} for asymmetric impacts; \texttt{abs\_peak\_deriv}, \texttt{wrms\_rot}
& Table~4 pairs RMS roll angular acceleration with road copying. Section~2.2 also permits P2P on angular body acceleration for isolated events. The two roll channels should be compared as alternative sources. \\

\texttt{Bounce\_rate\_6D}
& \texttt{rms\_deriv}, \texttt{p2p\_deriv}, \texttt{abs\_peak\_deriv}
& \texttt{wrms\_z\_d}, properly weighted transient metrics, raw-rate \texttt{rms}
& Differentiation produces heave acceleration. Table~4 directly pairs RMS sprung-mass vertical acceleration with vertical ride/bounce control, while Section~2.2 supports peak/P2P vertical acceleration for isolated obstacles. \\

\texttt{IMU\_LatAccelVal}
& none for a generic symmetric bump
& \texttt{p2p}, \texttt{rms}, \texttt{abs\_peak} for asymmetric impacts; \texttt{wrms\_xy} for cornering/repeated vibration
& The review associates lateral acceleration with choppiness and cornering, but Table~4 gives no direct generic-bump pairing. Its relevance depends on the event geometry. \\

\texttt{IMU\_YawRtVal}
& none for a generic bump
& raw-rate \texttt{rms}, \texttt{rms\_deriv}, \texttt{p2p\_deriv}, \texttt{wrms\_rot} for directional/asymmetric events
& Table~4 contains no direct yaw-to-bump or yaw-to-ride-attribute pairing. $W_e$ is mathematically valid for yaw angular acceleration, but predictive relevance is a handling/context hypothesis. \\
\bottomrule
\end{tabular}
\endgroup
\end{table*}

\FloatBarrier
\subsection{Evidence hierarchy and scope}

The review warns that the discussed KPIs are ``not universally applicable to every type of
manoeuvre and road profile'' \cite{guastadisegni2024}. Accordingly, evidence is treated in the
following order:
\begin{enumerate}
\item a metric paired with the same manoeuvre and subjective attribute in Table~4;
\item a metric explicitly recommended for the same physical response in Section~2.2;
\item a metric shown in the review's correlation examples, especially Figure~11;
\item a directionally or dimensionally valid extrapolation that still requires empirical testing.
\end{enumerate}
The review concerns attribute-specific ride assessment, whereas the current target is binary
feedback on short bump events. It therefore supplies a defensible prior over candidates, not a
guarantee that any one metric predicts the current label.

\subsection{Vertical impact and heave response}

Section~2.2 calls vertical absolute peak and P2P the ``most frequently selected indicators
\ldots{} during isolated events'' \cite{guastadisegni2024}. Table~4 makes the distinction more
specific: P2P seat vertical acceleration is paired with small and large impact on a bump,
whereas RMS seat vertical acceleration is paired with the corresponding impact on a cleat.
The same table pairs RMS sprung-mass vertical acceleration with vertical ride control or bounce
on bumpy roads with high wavelengths.

Consequently, \texttt{p2p} and \texttt{abs\_peak} have the strongest isolated-impact basis for
\texttt{IMU\_VerAccelVal}, while \texttt{rms} captures sustained body response and remains a
first candidate rather than a fallback. For \texttt{Bounce\_rate\_6D}, differentiation produces
the heave acceleration used by the paper. This makes \texttt{rms\_deriv} directly relevant to
bounce control, with \texttt{p2p\_deriv} and \texttt{abs\_peak\_deriv} capturing the transient
impact. The current centre-of-gravity measurements remain proxies for the seat and sprung-mass
locations used in the cited correlations.

\subsection{Pitch and roll body control}

Table~4 directly pairs RMS pitch angular acceleration with pitch control on flat and uneven
roads. It also lists settling time, P2P pitch angular acceleration, and RMS pitch angular
acceleration for pitch under acceleration or braking. The isolated-event discussion in
Section~2.2 states that P2P can also be evaluated on pitch and roll angular body acceleration.

For the rate channels, the matching quantities are obtained only after differentiation. Hence
\texttt{rms\_deriv} and \texttt{p2p\_deriv} are the first pitch candidates. For roll, Table~4's
direct evidence is RMS roll angular acceleration for road copying, making
\texttt{rms\_deriv} the first general candidate. \texttt{p2p\_deriv} becomes especially relevant
for a one-wheel or otherwise asymmetric bump. Absolute angular-acceleration peak and
$W_e$-weighted RMS are defensible additions, but the review does not establish them as the
primary correlation for these attributes.

\subsection{Longitudinal impact response}

Table~4 focuses its small- and large-impact rows on vertical seat acceleration. The direct
longitudinal evidence appears instead in Figure~11: front- and rear-seat-rail longitudinal P2P
are shown in correlation with impact harshness for a speed bump. This supports
\texttt{p2p} as the first longitudinal candidate. Absolute peak, RMS, and $W_d$-weighted RMS
remain plausible additions, but their bump-specific subjective correlation is less direct in this
review. Sensor-location differences again prevent interpreting the CG IMU feature as an exact
replica of the seat-rail KPI.

\subsection{Lateral and yaw responses}

The review relates lateral body acceleration to choppiness and discusses cornering manoeuvres,
while roll is used for road-copying and lane-change behaviour. It does not provide a Table~4
pairing between lateral acceleration or yaw and generic bump impact. These signals should
therefore have no unconditional first statistic for a nominally symmetric bump. They become
meaningful candidates when the data contain asymmetric road input, single-wheel impact,
cornering, or a directional correction after the obstacle. The condition is on the event, not
merely on the mathematical availability of a statistic.

\subsection{Why weighted and dose metrics are conditional}

The review lists WRMS as a ride-comfort KPI used for bumps as well as longer road profiles,
so it remains a legitimate candidate. However, its most direct bump-to-subjective-attribute
examples use P2P or RMS rather than WRMS. On a four-second event window, WRMS is also a
short-window frequency-shaped summary rather than a representative whole-body exposure
measurement.

RWRMS, MTVV, and VDV are motivated when ordinary RMS can underestimate occasional shocks
or strongly transient vibration. Their priority should therefore depend on a correctly computed
weighted crest factor or on a predefined transient-analysis design. The current raw
\texttt{crest}, \texttt{mtvv}, and \texttt{vdv} implementations can be explored as generic
shape and dose features, but they do not inherit the standards-based justification without
frequency weighting.

\subsection{Resulting interpretation}

For the present task, the paper supports a more precise statement than ``use peak statistics
for every signal'': P2P and absolute peak are first candidates for impact severity on the
vertical and longitudinal response, while RMS acceleration is a first candidate for heave,
pitch, and roll body control. P2P angular acceleration is first or near-first when the bump
induces a distinct pitch/roll transient. Lateral, yaw, frequency-weighted, and dose-based
features are conditional candidates whose value depends on event geometry, spectral content,
and the target driver's response.

\section{Unavailable or conditionally valid metrics}

\subsection{Third positive peak}
The review defines $z_{s,\mathrm{nod}}$ as the third positive peak of the \emph{sprung-mass
vertical displacement} $z_s(t)$ after an obstacle. Applying this operator directly to vertical
acceleration does not measure the same oscillation-decay property. It should be excluded unless
displacement is measured or reconstructed with a validated, drift-controlled estimator.

\subsection{Road-holding metrics}
The root integral of suspension-deflection PSD $\sigma_{SD}$, RMS suspension deflection,
suspension working space, and normalized dynamic wheel-load RMS require suspension travel
or vertical tyre-load signals. They must not be computed from body acceleration merely because
the same mathematical operation is available. Candidate raw channels for these quantities must
first be assigned verified physical meanings and units.

\section{Measurement and preprocessing limitations}

\begin{itemize}
\item ISO whole-body comfort evaluation is defined at body-contact locations such as the
seat surface, seat back, and feet. A 6DoF IMU near the vehicle centre of gravity is a useful
vehicle-response sensor, but its weighted features are proxies for occupant exposure.
\item Rotational acceleration must be expressed in $\mathrm{rad}/\mathrm{s}^2$ for physical
ISO reporting. A constant degree-to-radian conversion is immaterial after z-scoring, but it
is material for reported KPI values.
\item Per-axis WRMS values are not automatically comparable across directions. ISO total
vibration combines orthogonal components using direction-dependent multiplying factors.
Separate feature standardization intentionally removes this physical magnitude comparison.
\item The current 10-Hz smoothing cutoff suppresses much of the review's 10--25-Hz
secondary-ride/shake band. Cutoff and downsampling choices must be included in sensitivity
analysis if secondary ride is part of the target.
\item Short-window FFT weighting is affected by spectral resolution, window boundaries,
and leakage. The preprocessing and window duration are therefore part of the feature and
must remain identical across training and evaluation folds.
\end{itemize}

\section{Model-based feature selection principle}

Let $g=(\text{channel},\text{statistic})$ be the atomic candidate and let $S$ be a candidate
subset. Selection should be based on the hierarchical model's predictive objective, not on
the literature ordering or raw coefficient magnitude. For each held-out training driver $u$,
fit the population posterior on the remaining drivers and evaluate
\[
Q_t(S)=\frac{1}{U}\sum_u\frac{1}{|H_u|}
\sum_{i\in H_u}\log p(y_{ui}\mid\mathcal{D}_{-u},C_{u,t},S),
\]
where $t=0$ measures cold-start transfer and $t>0$ measures personalization after a fixed
context budget. The predictive importance of group $g$ is the paired difference
\[
\Delta_g=Q(S)-Q(S\setminus g).
\]
Channel-level importance removes all statistics belonging to a channel; statistic-level
importance removes one pair. Final selection should choose the smallest subset within one
standard error of the best macro-averaged held-out-driver score. The designated test driver
must remain untouched until this procedure is complete.

\section{Conclusion}

The review and standards determine which candidates are physically defensible, but they do
not determine the predictive winner for personalized binary feedback. For short bump-like
events, peak, P2P, and RMS acceleration features are first-class candidates; weighted RMS and
transient dose features provide complementary frequency and peak sensitivity. Feature
importance must ultimately be reported separately for population-common effects,
between-driver heterogeneity, and held-out-driver predictive gain.

\begin{thebibliography}{3}

\bibitem{guastadisegni2024}
G.~Guastadisegni, S.~De~Pinto, D.~Cancelli, \emph{et al.},
``Ride analysis tools for passenger cars: objective and subjective evaluation techniques
and correlation processes -- a review,'' \emph{Vehicle System Dynamics}, vol.~62, no.~7,
pp.~1876--1902, 2024, doi: 10.1080/00423114.2023.2259024.

\bibitem{iso2631}
International Organization for Standardization,
\emph{Mechanical vibration and shock -- Evaluation of human exposure to whole-body
vibration -- Part 1: General requirements}, ISO~2631-1:1997, 1997.

\bibitem{bs6841}
British Standards Institution,
\emph{Measurement and evaluation of human exposure to whole-body mechanical vibration
and repeated shock}, BS~6841:1987, 1987.

\end{thebibliography}

\end{document}
